\documentclass[../main.tex]{subfiles}
\begin{document}
%==========================================TIÊU ĐỀ TẬP========================================
\begin{table}[h]
    \begin{adjustwidth}{-1cm}{}
        \begin{tabular}{>{\centering\arraybackslash}p{6cm}|>{\raggedright\arraybackslash}p{12.5cm}}
            \multirow{5}{6cm}
            % Cột bên trái: ÉPISODE và số tập
            {\\[-40pt]
                \flaregothic\fontsize{20pt}{20pt}\selectfont \centering
                \color{eptype}ÉPISODE\color{black}\\
                \burbank\fontsize{72pt}{72pt}\selectfont
                \color{epnum}42\color{black}\\[6pt]
                \cafeteria\fontsize{20pt}{20pt}\selectfont\color{epnum}(3-12)\color{black}
                \\[18pt]
                \flaregothic\fontsize{14pt}{14pt}\selectfont
                \color{epattr}BIENVENUE, \uline{\color{Senchou} Marine} !\\
                \flaregothic\fontsize{18pt}{18pt}\selectfont
                \color{epattr}FIN DE TOME !\\[24pt]
                \color{black}
            }
             &
            % Dòng thứ nhất: tên tập tiếng Nhật
            {\fotskip\fontsize{18pt}{18pt}\selectfont コスプレ\ruby{女}{おんな}VS\ruby{緑}{みどり}\ruby{色}{いろ}のアレ
            }                                                                                      \\[6pt]
            % Dòng thứ hai: tên tập tiếng Anh
             & {\cafeteria\fontsize{18pt}{18pt}\selectfont \textit{hololive}'s snake}                                                                             \\[4pt]
            % Dòng thứ ba: tên tập tiếng Việt
             & {\vnmsans\fontsize{18pt}{18pt}\selectfont Rắn của ai?}                                                                                 \\[8pt]
            % Dòng thứ tư: Nhân vật xuất hiện (đã thêm Ryuuhou)
             & {\caviar{\fontsize{14pt}{14pt}\selectfont Nhân vật: \emp{kuon1} \emp{ryuuhou1} \emp{mel} \emp{fubuki} \emp{pekora} \emp{marine}}}
            \\[8pt]
            % Dòng cuối cùng: Thời gian diễn ra của tập
             & {\continuum\fontsize{16pt}{16pt}\selectfont Ngày phát hành: 23/02/2020}                                                                           \\[18pt]
        \end{tabular}
    \end{adjustwidth}
\end{table}
%=========================================TRUYỆN CHÍNH========================================
\color{black}

\txtbx[2]{{\color{vol} \uline{Tóm tắt:}} Một chú rắn lục bất ngờ xuất hiện trong văn phòng khiến Fubuki và Kuon hoảng loạn. Thủ phạm không ai khác chính là Thuyền trưởng Marine - người đã mang chú rắn từ trên núi về vì thấy nó\dots\ ``dễ thương.'' Một kế hoạch giải cứu ``ngoài không gian'' đầy vô lý của Senchou đã dẫn đến sự ra đời của một chòm sao mới và một cú đáp thảm họa xuống trường học của Kuon.}

Những ngày cuối cùng của mùa đông đang chầm chậm trôi qua, mang theo cái lạnh giá của tháng Hai len lỏi qua từng kẽ lá. Mặt trời lên muộn hơn, tỏa ra những tia nắng yếu ớt không đủ để xua tan cảm giác rét buốt còn sót lại trên những con phố Tokyo.

Bên trong văn phòng \hll, tiếng trò chuyện râm ran thường ngày bỗng chốc bị phá vỡ bởi một tiếng hét thất thanh vang dội khắp các phòng ban: ``\textbf{AAAAAAAAAAA!!!}''

Đó là tiếng hét đầy hoảng loạn của Shirakami Fubuki, khiến cả căn phòng như rơi vào trạng thái đóng băng. Kuon đang cặm cụi xử lý đống tài liệu tồn đọng, Mel đang sắp xếp lại mấy hộp quà từ người hâm mộ, và một cô nàng mới tới văn phòng lần đầu đang mải mê thử chiếc mũ cướp biển cho bộ đồ cosplay của mình. Cả ba lập tức quay phắt lại, đồng thanh hỏi với vẻ lo lắng: ``Có chuyện gì vậy, Fubuki-chan!?''

\chrint

Người mới xuất hiện đó chính là Houshou Marine (\ruby{宝}{ほう}\ruby{鐘}{しょう}マリン), một cô gái 17 tuổi có niềm đam mê mãnh liệt với việc hóa trang thành cướp biển. Dù hiện tại vẫn bị coi là một cosplayer, nhưng ước mơ cháy bỏng của Marine là tích góp đủ tiền từ công việc VTuber để mua một con tàu thật sự và trở thành thuyền trưởng hải tặc.

\model{24}{l}{10cm}{marine}

Marine nổi tiếng với tính cách mạnh dạn, ăn nói nhanh nhẹn nhưng đôi khi lại thiếu ý tứ, thường xuyên tung ra những câu đùa tục tĩu trêu chọc cả đồng nghiệp lẫn khán giả. Cô nàng hiếm khi kiềm chế bản năng hào hứng của mình, dẫn đến câu nói thương hiệu ``I'm Hornyyyy!\~♪'' đầy tai tiếng. Mặc dù tự nhận là nữ sinh 17 tuổi, nhưng giọng nói đậm chất ``bà cô'' và tầm hiểu biết sâu rộng về văn hóa đại chúng những năm 2000 khiến mọi người luôn nghi ngờ tuổi thật của cô phải tầm\dots\ 70.

Bù lại cho sự dâm đãng đôi khi quá đà, Marine lại là người có chỉ số giao tiếp (social skills) cực cao. Cô vô cùng nhạy cảm với bầu không khí xung quanh và luôn biết cách điều chỉnh hành vi để tránh gặp rắc rối trong môi trường công sở.

Bộ trang phục của Marine bao gồm một chiếc mũ cướp biển màu đen, một chiếc áo khoác đen (mà cô ấy thường mặc trễ vai), một chiếc áo khoác ngắn màu đỏ không tay có cúc màu vàng, viền vàng được trang trí bằng một chiếc ghim ve áo hình một con tàu cướp biển, mặc bên ngoài bộ đồ liền quần trong suốt có vòng cổ hoặc cổ áo xếp nếp màu đen, một chiếc dây nơ (ascot) đỏ có trâm cài, một chiếc váy ngắn xếp ly màu đỏ và trắng, một chiếc thắt lưng da màu nâu. Cô đeo găng tay trắng, đi quần tất dài đến đùi viền ren màu đen hoặc nâu và giày bốt cao cổ bằng da màu nâu có thể gập lại.

Cô có mái tóc nâu đỏ, được buộc thành hai bím thấp bằng hai chiếc ruy-băng màu đỏ. Cô có một đôi mắt hai màu: bên trái là mắt đỏ, và bên phải là mắt vàng (theo hướng của Marine), thường được che bằng một cái bịt mắt đen.

\chrint

Nàng Cáo trắng lúc này đang run rẩy kịch liệt, ngón tay chỉ về phía góc phòng, giọng nghẹn lại vì sợ hãi: ``Có rắn! Có rắn trong văn phòng kìa!''

Tất cả nhìn theo hướng tay cô và phát hiện một con rắn lục nhỏ đang trườn chậm rãi trên sàn nhà. Đôi mắt lạnh lùng và lớp vảy xanh mướt của nó khiến bầu không khí trong phòng chùng xuống vì bất an.

\img{3-13a}

Chứng kiến sinh vật không mời mà đến, Kuon khẽ giật mình, vội vàng lùi ra xa: ``Úi giời ôi, ở đâu ra cái giống này thế!?''

Mel bỗng vỗ tay lên trán như vừa sực nhớ ra điều gì đó: ``À, em có nghe ai đó nói là có người tìm thấy rắn ở ngọn núi quanh đây đấy\dots''

Fubuki quay phắt sang nhìn nàng Dơi, ánh mắt đầy tia máu vì hoảng loạn: ``\textbf{Là ai? Bữa giờ có ai lên núi chơi hả!?}''

Hai cô nàng lập tức đổ dồn ánh mắt về phía vị Tổng quản lý. Kuon vội vã xua tay thanh minh: ``Không phải anh nhé! Dạo này anh chỉ có đi học rồi tới đây làm việc, thời gian đâu mà leo núi? Với cả sức khỏe anh yếu xìu thế này thì leo trèo gì nổi\dots''

``\textbf{Thế thì ai!? Ai đã mang con rắn này xuống núi!?}'' Fubuki gào lên.

Ánh mắt của cả ba người đồng loạt chuyển sang Marine. Nàng Thuyền trưởng lúc này đang đứng im như phỗng, hai bàn tay khẽ đan vào nhau tạo thành một ám hiệu ngầm định rằng: ``Không phải tôi đâu nha.''

\img{3-13b}

Một sự im lặng bao trùm khắp văn phòng, chỉ còn tiếng sột soạt ghê người khi con rắn trườn trên sàn gỗ. Marine khẽ nuốt khan, nở một nụ cười gượng gạo: ``Này\dots\ không ai nghĩ là do tôi thật đấy chứ?''

Fubuki trừng mắt đầy nghi ngờ, còn Kuon thì thở dài ngao ngán, lẩm bẩm trong miệng: ``Chắc phải tìm cách xử lý con rắn này trước khi nó cắn ai đó đã\dots''

Tiếng rít nhẹ của con rắn vẫn vang lên đâu đó trong góc, nhưng giờ không còn ai để ý nữa. Tất cả sự chú ý đều tập trung vào Marine, người đang cố gắng đóng vai vô tội nhưng lại lộ rõ vẻ bối rối qua từng cử chỉ nhỏ.

Mel lên tiếng phá vỡ bầu không khí căng thẳng: ``Em cũng không nghĩ là Kuon-senpai đem nó về đâu.''

Kuon bật cười chua chát: ``Tất nhiên rồi. Người bình thường thấy rắn là chạy mất dép, trừ mấy ông gan lì đi săn rắn về ngâm rượu thì anh không nói.''

``Thật ạ?'' Mel ngạc nhiên hỏi.

Kuon gật đầu: ``Ừ, ở quê anh có những người như thế đấy. Nhưng số đó ít lắm.''

Fubuki gật gù, có vẻ đã bớt căng thẳng hơn một chút: ``Đ-Đúng rồi ha! Chẳng có ai điên rồ đến mức xách một con rắn độc về văn phòng cả\dots''

Thế nhưng, ánh mắt sắc lẹm của nàng Cáo trắng vẫn không rời khỏi nàng Thuyền trưởng - người lúc này đang ``diễn xuất'' vô cùng nhập tâm với đôi mắt nhắm nghiền, lưỡi thè ra và hai tay giơ lên trời như đang cầu nguyện sự vô tội.

\img{3-13c}

Kuon khẽ nghiêng người về phía Fubuki, thì thầm: ``Anh biết là ai làm rồi.''

Fubuki quay phắt lại, vẻ mặt đầy tò mò: ``Thế là ai? Nãy giờ chẳng thấy ai chịu nhận tội cả!''

Cậu Tổng nhún vai, rồi lững thững bước đến chỗ Marine với một sợi dây thừng chắc chắn trong tay. Marine giật mình lùi lại, giọng run rẩy nhưng vẫn không quên pha chút dâm đãng: ``Á\~ Kuon-senpai, anh định chơi trò gì với em vậy\dots?''

Kuon chẳng buồn đáp lại, anh chỉ buông một câu lạnh lùng: ``Ở yên đó một chút.''

\ct

Chỉ trong tích tắc, bằng những nút thắt điệu nghệ, Kuon đã trói chặt hai tay Marine ra sau lưng. Cô nàng đứng bất động, đôi mắt mở to nhìn cậu đầy sửng sốt.

``Xong,'' Kuon đáp gọn lỏn, phủi tay đầy mãn nguyện.

Thuyền trưởng nhíu mày: ``\dots\ Ể? Thế thôi ư?''

Cậu con trai không trả lời, anh quay lưng lại và nhanh chóng tiến về phía cửa. Fubuki đã đứng chờ sẵn ở đó với một nụ cười đầy ẩn ý, cô vẫy tay chào tạm biệt Marine: ``Vậy nhé, chúc em luôn cứng \uline{rắn} lên!''

Kuon dừng lại ở cửa, ngoái đầu nhìn nàng Cáo trắng rồi bật cười: ``Chơi chữ hay đấy Fubuki-chan. Chúc may mắn nhé, Thuyền trưởng!''

Dứt lời, cậu mở cửa bước ra ngoài cùng Fubuki, không quên đóng sầm cửa lại để mặc Marine trong căn phòng tối tăm cùng sinh vật xanh mướt kia.

Biết mình đã bị hai người họ bỏ rơi làm mồi cho rắn, Marine lập tức hét lên thảm thiết: ``\textbf{Khoan đã!} Em xin lỗi! Có gì chúng ta cứ bình tĩnh mà nói! Cởi trói cho em rồi mình thảo luận lại đi mà!''

\gif{3-13d}{1}{10}

Sau một hồi để Senchou van xin nài nỉ đến khản cả cổ, cuối cùng Kuon và Fubuki cũng chịu tha thứ. Mel khẽ thở dài, dùng kéo cắt bỏ sợi dây thừng đang siết chặt cổ tay Marine.

Nhưng vừa mới được tự do, Fubuki đã lập tức tra hỏi: ``Thế quái nào mà em lại xách con rắn này về đây hả!?''

Marine xoa xoa cổ tay bị đỏ ửng, ấp úng bào chữa: ``Chắc là do em bị\dots\ say núi ấy mà.''

Nghe lời giải thích phi lý đó, Kuon phì cười: ``Anh chỉ mới nghe say sóng, say xe, chứ chưa thấy ai bị say núi mà lại xách rắn về nhà bao giờ\dots\ Em quả là thiên tài sáng tạo đấy!''

Mel chớp mắt tò mò: ``Hải tặc như em thì leo núi làm gì thế Marine-chan?''

Kuon bồi thêm một câu đầy ẩn ý: ``Anh dám chắc là em lên đấy không phải để tìm kho báu đâu nhỉ?''

Marine giật mình, dường như bị nói trúng tim đen. Cô chưa kịp tìm cách biện minh thì Fubuki đã lao tới, nắm chặt vai cô và lắc mạnh liên tục: ``Senchou! Em phải có cách gì giải quyết mớ hỗn độn này đi chứ!? Đây là rắn lục đấy! Là rắn độc đấy em hiểu không!?''

Kuon đứng xem màn tra tấn mà cảm thấy như đang xem một vở kịch hài hạng sang. Anh khẽ thở dài: ``Thôi nào Cáo ơi\dots\ Bình tĩnh đi em, để người ta còn kịp thở mà nói chứ.''

Nàng Thuyền trưởng đầu óc quay cuồng vì bị lắc, cô hét lên: ``Khoan đã! Chị làm em chóng mặt quá rồi!''

Chỉ khi Kuon can thiệp, Fubuki mới chịu dừng lại, nhưng ánh mắt cô vẫn hằn học nhìn chằm chằm vào nàng Hải tặc.

``Vậy em nói đi! Tại sao em lại kéo con rắn này từ trên núi về đây!? Hay em đang âm mưu chiếm đoạt cái văn phòng này bằng vũ khí sinh học?'' Fubuki tiếp tục truy hỏi.

Marine vội vã xua tay: ``Ơ không không! Em không có âm mưu gì cả! Chỉ là\dots\ chỉ là em thấy nó dễ thương quá, nên em\dots\ lỡ tay mang nó về thôi mà\dots''

Cả ba người còn lại đứng hình mất vài giây. Fubuki há hốc mồm kinh ngạc: ``\textbf{Dễ thương!? Rắn lục mà em bảo dễ thương á!?}''

Kuon gập người cười sặc sụa, ôm bụng: ``Marine-chan ơi\dots\ Anh thề, gu thẩm mỹ của em đúng là 'khắm lọ' nhất cái văn phòng này luôn đấy!''

Mel nhẹ nhàng xác nhận lại: ``Marine-chan\dots\ Em thực sự thấy nó dễ thương sao?''

Senchou gật đầu lia lịa với vẻ mặt vô cùng nghiêm túc, khiến Fubuki càng thêm bất lực. ``Dù sao thì, mau đuổi nó ra khỏi đây đi!''

Marine lấy ra một món đồ chơi nhỏ và giải thích: ``Cái duy nhất em có lúc này là món đồ chơi Râu Đen mà em với Matsuri-senpai đã thức trắng đêm để chế tạo thôi.''

\img{3-13e}

Cả ba người nhìn cô bằng ánh mắt đầy hoài nghi. Marine vẫn cố gắng thuyết phục: ``Chắc là dùng nó làm đạn được mà! Chúng ta lấy tên Râu Đen này\dots\ phi thẳng vào con rắn!''

Fubuki nhìn Marine trân trối rồi hét lên: ``\textbf{Cái đó thì giúp ích được gì chứ!? Đây là rắn thật, không phải trò đùa đâu!}''

Marine cười khì khì hy vọng sẽ lay chuyển được mọi người, nhưng vô ích. Fubuki và Mel nhanh chóng quay sang bàn bạc phương án khác, bỏ mặc Marine đứng đó tự thân vận động.

Senchou lẩm bẩm tủi thân: ``Lại bị bơ luôn rồi kìa\dots''

Kuon nhìn cô với ánh mắt mệt mỏi: ``Tất nhiên là bị bơ. Ngoài cái ý tưởng điên rồ đó ra thì em còn cách nào khác không?''

Marine lắc đầu: ``Em chịu\dots\ Tay không làm sao bắt được rắn chứ\dots''

Kuon chép miệng thất vọng: ``Chán thật. Anh tưởng hải tặc là phải có súng ống đại bác gì chứ? Trên tàu thường có mấy thứ đó mà?''

Marine xụ mặt, chống hông phản bác: ``Nghèo rớt mùng tơi như em thì lấy đâu ra mấy thứ đó! Em chỉ là người thích hóa trang thành hải tặc thôi mà Kuon-senpai! Đừng có đặt kỳ vọng cao quá như thế!''

Cậu Tổng thở dài: ``Hết cách thật rồi\dots''

Fubuki vẫn đang đảo mắt quanh phòng, giọng run rẩy tìm kiếm sự trợ giúp: ``Có ai biết bắt rắn không!? Gọi cứu hộ hay ai đó đi! Em không chịu nổi thêm giây nào nữa đâu!''

Cáo trắng siết chặt vạt áo, lo sợ con rắn sẽ lao tới bất cứ lúc nào. Mel bỗng nảy ra một ý tưởng: ``Roboco-chan! Nếu là Roboco-chan thì chắc sẽ ổn đúng không?''

Fubuki sáng mắt lên như thấy ánh sáng cuối đường hầm: ``Ý hay đấy! Chị ấy là người máy, rắn có cắn cũng chẳng sao!''

Kuon nhíu mày đầy ngờ vực: ``Ý em là cái tên Robo-PON hậu đậu đó á? Anh thề là cái tên đó chỉ giỏi phá hoại thôi, có giúp được gì đâu.''

Nhưng cả Fubuki và Mel đều kiên quyết phản đối. Fubuki thốt lên: ``Nhưng dù sao cậu ấy cũng là người máy mà!'' Nàng Dơi còn bồi thêm: ``Có còn hơn không anh ạ!''

Kuon biết mình không thể thắng được hai cô nàng này, đành phẩy tay đồng ý: ``Thôi được rồi, mấy đứa thử đi gọi đi. Để xem có cơm cháo gì không.'' Anh thầm nhủ: ``Bó tay toàn tập với hội này.''

Được lệnh, Fubuki hào hứng chạy vụt ra ngoài: ``Để em đi tìm chị ấy ngay!''

\ct

Năm phút sau, nàng Cáo trắng trở lại. Thế nhưng, thay vì Roboco, cô lại lôi theo một nàng Thỏ tóc xanh đang gà gật ngủ gật. Fubuki thở hồng hộc, mồ hôi nhễ nhại phân trần: ``E-Em không tìm thấy Roboco-san đâu cả, nên đành hốt đại người này về đây!''

\img{3-13f}

Mel nhìn thấy ``vị cứu tinh'' mới thì ôm đầu thất vọng: ``Trời ơi, sao lại vớ ngay cái đứa dở nhất trong vụ này về thế này\dots''

Kuon cười chua chát: ``Thôi xong, xôi hỏng bỏng không thật rồi. Ai chẳng biết thỏ là món khoái khẩu của rắn chứ.''

Nghe thấy tên mình, Pekora lờ mờ mở mắt, ngơ ngác nhìn quanh. Mel tiến tới, cố gắng vớt vát chút hy vọng cuối cùng: ``Pekora-chan\dots\ Em có thể xử lý nó giúp mọi người được không?''

Ngay khi Mel chỉ tay về phía con rắn đang trườn trên sàn, nàng Thỏ sợ đến mức mặt không còn giọt máu, đầu lắc như điện giật. Không chịu nổi áp lực kinh hoàng, Pekora lăn đùng ra bất tỉnh nhân sự ngay tại chỗ.

\img{3-13g}

Mel thảng thốt la lên: ``Em ấy xỉu luôn rồi kìa!''

Kuon đứng cạnh chỉ biết lắc đầu ngao ngán: ``Đã bảo là Fubu-chan tìm sai người mà lại.''

Mel gật đầu đồng ý: ``Đúng thế! Mau đưa em ấy ra ngoài ngay đi!''

Nói rồi, Kuon nhanh chóng bế nàng Thỏ tội nghiệp ra khỏi phòng, để lại hai cô gái tiếp tục đối mặt với tình huống tréo ngoe.

% -------------------- BỔ SUNG RYUUSHOU --------------------
\ct
Vừa lúc Kuon đặt Pekora lên ghế ngoài hành lang, một giọng nói lanh lảnh vang lên từ cuối dãy:
``Onii-sama làm gì mà nhộn nhịp thế? Lôi cả xác thỏ ra đây là sắp làm thịt à?''

Kuon ngước lên, thấy cô em gái Ryuuhou đang chống nạnh, mái tóc dài xõa lưng, đôi mắt hai màu (phải đỏ, trái xanh lam) híp lại đầy tinh nghịch. Hôm nay em mặc một chiếc áo phông trắng và quần soóc jean, trông như vừa đi đâu về.

``Ryuuhou? Em đến đây làm gì?'' Kuon ngạc nhiên.

``Thì em sang chơi với Subaru-senpai, nhưng nghe thấy tiếng hét ầm ĩ quá nên chạy ra xem. Có rắn à?'' Ryuuhou nhón chân nhìn qua cửa, mặt vô cùng phấn khích.

``Rắn lục đấy, em tránh xa ra.'' Kuon vội kéo em lại.

Nhưng Ryuuhou đã nhanh chân luồn vào phòng, mắt sáng rỡ: ``Oa, xanh đẹp quá! Màu này hợp làm phụ kiện cho nhóm nhạc của em ghê!''

Fubuki và Mel đang co rúm ở góc tường, thấy con bé lao vào thì hét lên: ``Ra ngoài đi Ryuuhou-chan! Nguy hiểm!''

Ryuuhou cười khúc khích: ``Có gì đâu, rắn nhỏ thế này mà. Ở quê ngoại em còn thấy rắn hổ mang to như cột đình cơ.''

Cô bé chậm rãi tiến lại gần con rắn, khom người. Con rắn lục bỗng dựng cổ lên, lưỡi thè ra tít tít. Ryuuhou bất ngờ chụp nhanh sau gáy nó, nhấc bổng lên với vẻ mặt đầy tự hào: ``Đây nhé! Bắt rắn không khó đâu!''

Cả phòng chết lặng. Marine há hốc mồm: ``Em\dots\ em bắt được rồi sao?''

Ryuuhou vừa giữ chặt miệng rắn vừa quay sang nhìn Kuon, cười đểu: ``Onii-sama đúng là vô dụng. Làm quản lý mà sợ rắn à?''

Kuon mặt đỏ tía tai: ``Ai bảo anh sợ! Anh chỉ không muốn mạo hiểm thôi!''

Fubuki chạy lại, mắt long lanh: ``Ryuuhou-chan em giỏi quá! Mau bỏ nó vào bao mang đi!''

Ryuuhou liệc nhìn con rắn, rồi bỗng ném nó về phía Marine: ``Trả lại cho chủ nhân của nó đi ạ!''

Con rắn văng trúng ngực Marine, cô nàng hét thất thanh, nhảy dựng lên, rồi vô tình đạp phải thùng đồ chơi Râu Đen. Cánh cửa bệ phóng bật mở, hút Marine và con rắn vào trong một lần nữa.

``Á á á! Lại nữa rồi!'' Marine gào lên khi chiếc thùng lao vút lên trần nhà.

Ryuuhou đứng nhìn, tay che miệng cười: ``Em chỉ đùa tí thôi mà, sao nó nghiêm trọng thế?''

Kuon ôm đầu bất lực: ``Em đúng là 'lợi hại' hơn anh nghĩ nhiều!''

\ct
% -------------------- KẾT THÚC BỔ SUNG --------------------

Chứng kiến cảnh tượng hỗn loạn, Marine (lần thứ hai) hét lên từ bên trong thùng: ``Mọi người tránh ra! Để em xử lý!''

Fubuki tròn mắt nhìn chiếc thùng: ``\textbf{Cái quái gì thế này!?}''

Nàng Hải tặc đặt tay lên thùng rượu, dõng dạc tuyên bố: ``Em sẽ dùng món bảo bối Râu Đen siêu to khổng lồ này để tống khứ con rắn ra ngoài không gian!''

Ý tưởng điên rồ của Senchou khiến cả Fubuki lẫn Kuon đều há hốc mồm. Họ đồng thanh thốt lên trong sự ngỡ ngàng: ``\textbf{CÁI GÌ CƠ!?}'' Mel cũng đứng hình, không nói nên lời. Ryuuhou thì vỗ tay thích thú: ``Hay quá! Em muốn xem!''

Không để ai kịp phản ứng, Marine nhanh chóng tóm lấy con rắn (lần nữa), nhảy tót vào trong bệ phóng cùng nó và hét lớn: ``\textbf{Nào! Mày sẽ phải cùng tao lên đường thôi!}''

Kuon hốt hoảng giơ tay ngăn cản: ``Ơ từ từ đã, có nhất thiết em phải chui vào trong đó cùng nó không-``

Lời Kuon chưa dứt thì từ hư không bỗng xuất hiện hàng tá thanh kiếm nhựa bay tới, cắm phập vào chiếc thùng rượu khổng lồ. Một phản lực cực mạnh bùng phát, phóng Marine cùng con rắn vọt lên cao như một quả tên lửa thực thụ.

\begin{figure}[H]
  \centering
  \begin{subfigure}[t]{0.45\textwidth}
    \centering
    \animategraphics[autoplay,loop,width=8cm]{12}{../../../../Image/Book3/3-13i/3-13i.}{1}{51}
    \caption{Hàng loạt kiếm nhựa xuyên vào cái thùng đồ chơi\dots}
  \end{subfigure}
  \quad
  \begin{subfigure}[t]{0.45\textwidth}
    \centering
    \animategraphics[autoplay,loop,width=8cm]{12}{../../../../Image/Book3/3-13i/3-13i.}{52}{108}
    \caption{\dots biến Marine thành một tên lửa bay thẳng vào không gian.}
  \end{subfigure}
\end{figure}

Âm thanh chói tai vang lên khi ``tên lửa'' đâm thủng cả trần nhà, khiến Fubuki thét lên trong sự kinh hoàng: ``Thuyền trưởng ơi\dots!''

Mel và Kuon ngẩn ngơ nhìn theo lỗ hổng trên trần nhà. Một lúc sau, Kuon mới thốt lên được một câu: ``Vãi thật\dots\ Thế cũng làm được luôn à? Nhưng mà anh cứ thấy có gì đó sai sai\dots''

Mel gật đầu, mặt đầy vẻ bối rối: ``Em cũng thấy thế\dots\ Nếu em ấy đã bắt được con rắn rồi, sao không mang nó trả về núi cho xong?''

Kuon tặc lưỡi, nhún vai thở dài: ``Thôi thì cứ cho là logic của Thuyền trưởng nó khác người thường đi em. Anh bó tay thật rồi.''

Ryuuhou từ nãy giờ vẫn đứng khoanh tay nhìn, bỗng lên tiếng: ``Em thấy vui vãi. Onii-sama, mai em lại qua chơi nhé!''

Kuon quắc mắt: ``Đừng có mà phá nữa!''

Họ đứng đó, lặng lẽ dõi mắt nhìn theo vệt khói của ``tên lửa Marine'' đang nhạt dần trên bầu trời, mang theo cả con rắn và những ý tưởng điên rồ nhất thế gian. Ryuuhou thì thầm vào tai Kuon: ``Em cá là chị ấy sẽ rơi đúng vào trường học của anh đấy.''

``Im đi, em gái lắm mồm!'' Kuon đáp, nhưng trong lòng thầm nghĩ biết đâu em nó nói đúng.

%==========================================TRUYỆN PHỤ=========================================
\omk

\begin{center}
  \pov{Hikawa Kuon}
  \textit{\uline{Tại sân thượng trường Đại học Xây dựng Kawachi.}}
\end{center}

Buổi tối hôm đó tại trường, bầu trời đêm vô cùng quang đãng với ánh trăng dịu nhẹ chiếu rọi xuống sân trường. Tôi cùng hai người bạn thân là Yamazu và Miharu đã rủ nhau lên tầng thượng của dãy nhà H1 để ngắm sao và hoàn thành bài tập chiêm tinh cho môn lập trình. Nhiệm vụ của chúng tôi là thu thập dữ liệu về các chòm sao rồi viết một chương trình phân tích quỹ đạo của chúng.

Để buổi học thêm phần chuyên nghiệp, Yamazu đã vác theo một chiếc kính viễn vọng loại xịn. Miharu thì lo phần hậu cần với máy phát điện di động và một hộp đồ ăn nhẹ. Cả ba chúng tôi ngồi quây quần bên máy tính, hăng hái bắt tay vào việc.

Không khí trên tầng thượng vô cùng yên bình, chỉ có tiếng gió thổi vi vu và tiếng quạt tản nhiệt của laptop hoạt động khe khẽ. Chúng tôi phân công rõ ràng: Yamazu và Miharu thay phiên nhau quan sát, còn tôi chịu trách nhiệm tổng hợp báo cáo.

Mọi chuyện đang diễn ra hết sức thuận lợi thì Miharu đột nhiên lên tiếng bằng giọng đầy vẻ nghi hoặc: ``Ê Kuon, tôi thấy có chòm sao này lạ lắm ông ạ!''

Tôi đang mải gõ code, ngẩng đầu lên hỏi: ``Gì thế?''

``Tôi chưa bao giờ thấy chòm sao này trong sách giáo khoa cả. Tra Google cũng chẳng thấy kết quả nào trùng khớp luôn!''

Tò mò, cả ba chúng tôi cùng ghé mắt vào kính viễn vọng. Yamazu, vốn là người điềm tĩnh nhưng lúc này cũng phải đùa cợt: ``Trông nó cứ như\dots\ con gì lai giữa con rắn với cái cục gì tròn tròn ấy nhỉ?''

\img{3-13j}

Tôi nhíu mày nhìn kỹ hơn. Chỉ một thoáng sau, hình ảnh cái thùng rượu khổng lồ lúc sáng hiện ra trong đầu, tôi không nhịn được mà bật cười: ``Quái lạ thật\dots\ À, tôi biết nó là cái gì rồi!''

``Biết gì cơ?'' Yamazu hỏi dồn.

Tôi chỉ vào chòm sao lạ hoắc kia, dõng dạc tuyên bố: ``Đó chính là chòm sao Tỵ!''

Cả hai cậu bạn ngơ ngác nhìn nhau: ``Chòm sao Tỵ? Ý ông là chòm sao con Rắn á?''

Tôi nén cười, giải thích cho đúng 'nguồn gốc' mà mình vừa chứng kiến: ``Thấy cái phần mảnh dài ở đầu không? Nó giống hình con rắn lục lúc sáng\dots\ à không, giống con rắn. Còn cái cục to đùng phía sau\dots\ người ta bảo đó là cái thùng rượu đấy!''

Miharu nghiêng đầu khó hiểu: ``Thùng rượu gì cơ?''

``Cái này thì tôi chịu. Chỉ biết là thiên hạ đặt tên như thế nên tôi gọi theo thôi!'' Câu nói của tôi khiến cả đám bạn phá lên cười trước phát hiện kỳ thú này.

Dù vậy, Yamazu vẫn không khỏi thắc mắc: ``Nhưng mà lạ thật, sao chòm này lại không có trong tài liệu chính thống? Hay là chúng ta vừa mới phát hiện ra một thiên thể mới?''

Miharu lập tức dùng điện thoại chụp lại hình ảnh qua kính viễn vọng, hớn hở: ``Chắc chắn phải điều tra kỹ hơn mới được! Biết đâu tên ba đứa mình lại được ghi danh vào lịch sử thiên văn học ấy chứ!''

Tôi phì cười, chống cằm nhìn lên trời cao: ``Đừng có mơ mộng hão huyền. Có khi đây chỉ là trò nghịch ngợm của \uline{ai đó} thôi, chứ chòm sao thật làm gì có hình thù quái đản thế kia.''

Dù nói vậy, nhưng sự tò mò trong lòng cả ba vẫn không hề thuyên giảm, chúng tôi quyết định dành thêm thời gian để nghiên cứu ``chòm sao Tỵ'' đầy bí ẩn này.

\ct

Sau một hồi loay hoay với chòm sao mới, bầu trời bỗng rực sáng bởi vệt sáng kéo dài của một ngôi sao băng. Cả đám lập tức hào hứng hẳn lên.

Yamazu là người đầu tiên chắp tay, nhắm mắt lầm rầm: ``Tôi ước bài kiểm tra Vật lý đại cương lần này được qua môn!''

Miharu cũng nhanh chóng hùa theo: ``Mong cho bài tập nhóm này được điểm A+!''

Tôi ban đầu chỉ ngồi cười nhìn hai đứa bạn, cũng định ước điều gì đó. Nhưng khi ngước nhìn kỹ vệt sáng kia, tôi đột nhiên khựng lại, nét mặt dần trở nên nghiêm trọng. Tôi khẽ nói, giọng run nhẹ: ``Ờm\dots\ Mọi người ơi?''

``G-Gì vậy?'' Yamazu đáp.

``Cái này\dots\ không giống sao băng lắm.'' Tôi chỉ tay lên trời.

Miharu ngước nhìn, mặt cắt không còn hột máu: ``\dots Hả?''

Chúng tôi cùng dõi mắt theo vầng sáng đang lớn dần với tốc độ kinh hồn. Tôi nhận ra có điều gì đó cực kỳ không ổn. Vệt sáng này không hề rơi theo quỹ đạo tự nhiên mà đang\dots\ lao thẳng về phía tầng thượng chúng tôi đang ngồi!

Tôi hoảng hốt hét lên: ``\textbf{MAU RA KHỎI ĐÂY NGAY!!!}''

Yamazu và Miharu giật mình: ``Gì cơ!?'' ``Ông đùa tụi tôi à!?''

Không đợi hai đứa bạn kịp phản ứng, tôi túm lấy áo cả hai rồi kéo mạnh vào góc khuất sau bức tường ngăn bể nước.

*RẦM!!!*

Một tiếng nổ kinh thiên động địa vang lên khiến cả dãy nhà H1 rung chuyển dữ dội, bụi bặm bay mù mịt. Một lúc sau, khi tiếng động đã dứt, ba đứa mới dám thò đầu ra nhìn.

Trước mắt chúng tôi là một cảnh tượng không thể tin nổi: một cái hố lớn xuất hiện ngay giữa sân thượng, và ngay trung tâm cái hố đó là một cô gái với mái tóc nâu đỏ, hai bím tóc thấp, đang lồm cồm bò dậy với vẻ mặt đầy đau đớn.

Tôi dẫn đầu đám bạn tiến lại gần. Khi nhận ra người vừa ``hạ cánh'' là ai, tôi chỉ biết ngửa mặt lên trời mà cười: ``Vãi cả Senchou\dots\ chào buổi tối nhé. Chuyến du lịch không gian về có vui không?''

Marine vừa xoa đầu vừa rên rỉ: ``Đ-Đau quá\dots\ Không ngờ là mình lại hạ cánh trượt mục tiêu xa đến thế này.''

Hai cậu bạn của tôi đứng đực ra như phỗng trước cô gái lạ mặt. Yamazu lắp bắp: ``Đ-Đ-Đi du lịch vũ trụ!? C-Cô nói cái quái gì thế!?'' còn Miharu cũng hỏi: ``Mà khoan, cô nàng này là ai vậy Kuon?''

Tôi thở dài, quay sang nhìn hai người bạn đang tái mét mặt mày: ``Chuyện dài lắm, thôi thì nhân tiện Marine-chan ở đây, để chúng tôi kể cho hai ông nghe về một ngày 'bình thường' tại văn phòng \hll\dots''

Đêm đó, chúng tôi ngồi quây thành vòng tròn ngay trên tầng thượng. Marine hào hứng kể lại hành trình phi logic của mình: từ việc bị phóng lên không gian cho đến lúc cô nhầm con rắn là ``kho báu vũ trụ'' nên quyết định xách nó quay về Trái đất.

Cậu bạn lập trình trợn tròn mắt: ``Cái gì cơ? Cô mang con rắn đó từ ngoài vũ trụ về á!?''

Yamazu thì lắc đầu ngao ngán: ``Đây đúng là chuyện điên rồ nhất mà tôi từng được nghe trong đời.''

Tôi chỉ biết cười khổ: ``Mấy cô nàng ở chỗ tôi làm đều như thế cả đấy. Hai ông rồi cũng sẽ quen thôi.''

Nàng Hải tặc vẫn say sưa kể về những điều kỳ diệu cô thấy trên quãng đường trở về. Khi nghe chúng tôi nhắc đến ``chòm sao Tỵ'' mới phát hiện, Marine vỗ đùi bôm bốp: ``Đúng rồi! Chính là em đấy! Đó là vệt khói mà em để lại khi bay ngang qua các ngôi sao đó!''

Yamazu và Miharu chỉ biết nhìn nhau trân trối, còn tôi thì không nhịn được cười: ``Ra là vậy. Anh đúng là không ngờ tới cái 'sáng chế' này luôn đấy.''

Tôi hối thúc hai người bạn: ``Nhớ ghi vào báo cáo nhé: 'Chòm sao Rắn Thùng' là công trình của Thuyền trưởng Houshou Marine!''

Cuối cùng, nhìn cái hố to tướng trên trần nhà trường học, tôi chỉ biết tặc lưỡi tự nhủ: ``Buổi làm việc đầu tiên của em mà đã phá hoại thế này rồi\dots\ thôi thì coi như anh tha cho lần này.''

Tôi lặng lẽ ngồi bên cạnh, nhìn Marine đang cười đùa vui vẻ kể chuyện cho hai cậu bạn của mình nghe. Tôi không thể không nghĩ về sự hỗn loạn mà cô nàng này mang lại, từ văn phòng cho đến tận trường học của tôi.

``\textit{Senchou ơi\dots\ Mở màn mà đã hoành tráng đến mức này thì\dots\ siêu thật đấy. Anh nghĩ em nên chuyển nghề làm không tặc luôn đi cho bõ công phá hoại.}''

Tôi mỉm cười, lòng cảm thấy nhẹ nhõm lạ kỳ. Dù chẳng bao giờ hiểu hết được logic của các cô gái ở công ty, nhưng ít ra cuộc sống của tôi đã bớt tẻ nhạt hơn rất nhiều kể từ khi có họ.

``\textit{Rồi không biết mình sẽ còn bị cuốn vào những điều điên rồ nào nữa đây\dots}'' tôi nghĩ thầm, trong lòng thầm hứa sẽ luôn sẵn sàng đối mặt với những câu chuyện ``mất não'' tiếp theo của đại gia đình \hll.

\end{document}